\documentclass[11pt]{article}
\usepackage[margin=1.2in]{geometry}
\usepackage{booktabs,enumitem,hyperref}
\usepackage{fontspec}
\setmainfont{TeX Gyre Pagella}

\title{The Ledger of Meluhha --- Reproducibility Artifact}
\author{Rajeshkumar Venugopal\\
  \small Third Buyer Advisory LLC, Waterford, Michigan\\
  \small\href{mailto:vrajeshkumar@gmail.com}{vrajeshkumar@gmail.com}}
\date{April 2026}

\begin{document}
\maketitle

\section{What this is}

Reproducibility artifact for the working paper
\textit{The Ledger of Meluhha --- Indus Valley Script as
Metrological Accounting Code} (Venugopal 2026).

Every claim in the paper is backed by data in the SQLite
databases below.
Every computation is reproducible with the F\# scripts below.
No proprietary tools, no cloud services, no API keys.

\section{Compile the paper (easiest: Overleaf)}

If you only want to read, edit, or recompile the paper,
Overleaf is the fastest path. No local install required.

\begin{enumerate}
\item Go to \url{https://www.overleaf.com} (free account).
\item \textbf{New Project} $\to$ \textbf{Upload Project} $\to$
      upload the zip of this artifact (or just
      \texttt{indus\_accounting.tex} + \texttt{citations.lua}).
\item \textbf{Menu} (top-left) $\to$ \textbf{Compiler} $\to$
      select \textbf{LuaLaTeX}.
      This is required --- the paper uses \texttt{fontspec},
      \texttt{luacode}, and \texttt{\textbackslash directlua}
      which only work under LuaLaTeX.
\item Click \textbf{Recompile}. Two passes resolve
      cross-references automatically.
\end{enumerate}

\noindent
Overleaf's free tier includes LuaLaTeX, TeX Gyre Pagella,
and every package this paper uses. No credit card, no install.

\section{Full reproducibility (local)}

To run the F\# scripts and rebuild the databases locally:

\begin{itemize}
\item \texttt{sqlite3} --- query any database directly
\item \texttt{dotnet fsi} (F\# 8+ / .NET 8+) --- run the scripts
\item \texttt{lualatex} (TeX Live 2024+) --- compile the paper locally
\item \texttt{java} (17+) + \texttt{alloy.jar} --- verify the schema
\item \texttt{cmake} (3.20+) --- build system
\end{itemize}

\noindent
macOS: \texttt{brew install dotnet texlive sqlite openjdk@17 cmake}\\
Linux: \texttt{sudo apt install dotnet-sdk-8.0 sqlite3 texlive-full default-jdk cmake}\\
Windows: not supported --- use WSL~2 (Ubuntu). See \texttt{CMakeLists.txt}.

\section{Files}

\begin{description}[style=nextline,leftmargin=2em]

\item[\texttt{indus\_accounting.tex}]
  The paper. LuaLaTeX source. Tables generated via embedded Lua.
  Compile with \texttt{lualatex -shell-escape indus\_accounting.tex}
  (two passes for cross-references).

\item[\texttt{indus\_accounting.pdf}]
  Compiled paper (26 pages).

\item[\texttt{citations.lua}]
  Lua citation engine loaded by the paper. Implements first-use
  full citation, subsequent \textit{op.\ cit.}

\item[\texttt{indus\_codebook.db}]
  SQLite. The metrological codebook: 11 weight tiers, 8 commodity
  codes, 5 routes, 6 quantity codes, 5 merchant marks, 28
  sign-to-role mappings, 3 positional fallback rules.
  Created by \texttt{indus\_seed\_codebook.fsx}.

\item[\texttt{indus\_corpus.db}]
  SQLite. Normalised corpus database: 2{,}373 signs from 3 sources
  (Rajan--Sivanantham 42, Tamil Treebank 1{,}929, Mahadevan 402),
  7{,}644 sign occurrences, 600 inscriptions, 18 sites, 20
  morphological parallels. Schema verified by Alloy (9/9 UNSAT).
  Created by \texttt{indus\_ingest\_corpus.fsx}.

\item[\texttt{indus\_lssc.db}]
  SQLite. Transition entropy and LSSC scores per sign and per
  inscription for the Tamil proxy corpus.
  Created by \texttt{indus\_lssc.fsx}.

\item[\texttt{indus\_seed\_codebook.fsx}]
  F\# script. Seeds \texttt{indus\_codebook.db} from the paper's
  schema. Usage: \texttt{dotnet fsi indus\_seed\_codebook.fsx codebook.db}

\item[\texttt{indus\_decoder.fsx}]
  F\# script. Cargo-tag decoder. Reads all configuration from
  \texttt{indus\_codebook.db}. Zero hardcoded values.
  Usage: \texttt{dotnet fsi indus\_decoder.fsx codebook.db}

\item[\texttt{indus\_ingest\_corpus.fsx}]
  F\# script. Builds the corpus database from Tamil Treebank CSVs,
  Rajan--Sivanantham base sign table, and Mahadevan histogram signs.
  Usage: \texttt{dotnet fsi indus\_ingest\_corpus.fsx <indus-valley-dir> corpus.db}

\item[\texttt{indus\_lssc.fsx}]
  F\# script. Computes transition entropy, closure/option
  classification, LSSC scores, and TSF Condition~2 check.
  Usage: \texttt{dotnet fsi indus\_lssc.fsx <corpus.csv> lssc.db}

\item[\texttt{indus\_tn\_decoder.fsx}]
  F\# script. Cross-corpus Tamil Nadu potsherd decoder. Resolves
  TN graffiti through IVC codebook via morphological parallels.
  Usage: \texttt{dotnet fsi indus\_tn\_decoder.fsx corpus.db codebook.db}

\item[\texttt{spec/indus\_corpus.als}]
  Alloy 6 specification. Nine structural assertions on the corpus
  schema. All UNSAT at scope~6 (zero counterexamples).
  Run: \texttt{java -jar alloy.jar exec spec/indus\_corpus.als}

\item[\texttt{data/morphological-parallels/rs2025\_plate\_parallels.csv}]
  CSV. Twenty morphological parallels transcribed from
  Rajan--Sivanantham (2025) comparison plates pp.~6--7.
  Each row cites the specific plate, panel, and row position.

\end{description}

\section{Quick start (CMake)}

\begin{verbatim}
# Configure (point to the indus-valley data directory)
cmake -S . -B build -DINDUS_VALLEY_DIR=../indus-valley

# Build everything: codebook + corpus + LSSC
cmake --build build

# Individual targets:
cmake --build build --target codebook    # seed codebook DB
cmake --build build --target corpus      # ingest corpus + parallels
cmake --build build --target lssc        # LSSC transition entropy
cmake --build build --target decode      # run IVC cargo-tag decoder
cmake --build build --target tn-decode   # TN cross-corpus decoder
cmake --build build --target verify      # Alloy schema verification
cmake --build build --target paper       # compile the paper (lualatex)
cmake --build build --target readme      # compile the README
cmake --build build --target clean-db    # remove generated DBs
\end{verbatim}

\section{Query the results directly}

\begin{verbatim}
sqlite3 indus_codebook.db "SELECT * FROM sign_role WHERE role='commodity'"
sqlite3 indus_corpus.db "SELECT * FROM morphological_parallel"
sqlite3 indus_lssc.db "SELECT class, COUNT(*), ROUND(AVG(entropy),4)
                        FROM sign_entropy GROUP BY class"
\end{verbatim}

\section{License}

CC~BY~4.0. Data from Rajan and Sivanantham (2025, 2026) used
under fair academic use for morphological comparison.
Tamil Treebank v0.1 under CC~BY-NC-SA~3.0 (Ramasamy \&
Abokrtsk 2011).

\section{Citation}

\begin{verbatim}
Venugopal, R. (2026). The Ledger of Meluhha: Indus Valley
Script as Metrological Accounting Code. Working Paper.
Third Buyer Advisory LLC. Zenodo.
doi: [to be assigned]
\end{verbatim}

\end{document}
